\documentclass[10pt,a4paper]{amsart}
\usepackage[margin=19mm]{geometry}
\usepackage{amsmath,amssymb,mathtools}
\usepackage[scr=rsfs,cal=euler]{mathalfa}
\usepackage{xfrac}
\usepackage[unicode,hidelinks,bookmarks=false]{hyperref}
\newcommand{\ie}{i.e.\ }
\newcommand{\cf}{cf.\ }
\newcommand{\resp}{respectively\ }
\newcommand{\Subsection}{\subsection}
\providecommand{\nobreakdash}{\nobreak\hbox{-}}
\setlength{\emergencystretch}{2em}
\multlinegap0pt
\usepackage{graphicx}
\usepackage{enumerate}

\newcommand{\darkblue}{}
\newcommand{\darkred}{}
\newcommand{\lightgrey}{}
\newcommand{\defn}[1]{\emph{#1}}
\newcommand{\chk}[1]{#1^{\normalfont\smash{\scalebox{.5}[.5]{\rotatebox{90}{\textbf{\textless}}}}}}

\newcommand{\Z}{\mathbb{Z}}
\newcommand{\R}{\mathbb{R}}
\newcommand{\lcm}{{\sf lcm}}
\newcommand{\Expt}[2]{\operatorname*{\mathbb{E}}\limits_{#1}(#2)}
\newcommand{\Var}[2]{\operatorname*{\mathbb{V}}\limits_{#1}(#2)}
\newcommand{\maximum}[2]{\operatorname*{max}\limits_{#1}(#2)}
\newcommand{\modn}{\text{ }\mathrm{mod}\text{ }n}

\newcommand{\wa}{\widetilde{W}}
\newcommand{\wphi}{\widetilde{\Phi}}
\newcommand{\h}{\mathrm{ht}}
\newcommand{\Q}{\mathcal{Q}}
\newcommand{\ac}{\mathcal{A}}
\newcommand{\cw}{\chk{\omega}}
\newcommand{\Sommers}{\mathcal{S}}
\newcommand{\bijQtoW}{\mathsf{crt}}

\newcommand{\core}{{\sf core}}
\newcommand{\cwl}{\chk{\Lambda}}
\newcommand{\wex}{\wa_{\mathrm{ex}}}
\newcommand{\cycl}{\Omega}
\newcommand{\waf}{\widetilde{w}}
\newcommand{\wf}{\waf_{\rat}}
\newcommand{\amax}{\tilde{\alpha}}
\newcommand{\rhoh}{\frac{\chk{\rho}}{h}}

\newcommand{\size}{{\sf size}}
\newcommand{\sizeb}{\size^{(b)}}
\newcommand{\inv}{{\sf inv}}

\newcommand{\w}{\mathsf{w}}
\newcommand{\x}{\mathbf{x}}
\newcommand{\rat}{b}
\newcommand{\bt}{t_b}
\newcommand{\br}{r_b}

\newtheorem{theorem}{Theorem}[section]
\newtheorem{prop}[theorem]{Proposition}
\newtheorem{lemma}[theorem]{Lemma}
\newtheorem{coro}[theorem]{Corollary}
\theoremstyle{definition}
\newtheorem{defi}[theorem]{Definition}
\theoremstyle{remark}
\newtheorem{rema}[theorem]{Remark}
\newtheorem{exam}[theorem]{Example}

\usepackage{cleveref}

\title[Expected core size, types B and C]{Strange expectations in affine Weyl groups: original Theorem 7.1, types $B$ and $C$}
\author{Eric Nathan Stucky, Marko Thiel and Nathan Williams}
\date{Source: Algebraic Combinatorics 7 (2024), publisher version; DOI 10.5802/alco.383}
\begin{document}\maketitle
\noindent\textit{Editorial scope.} The counted unit is the original printed Theorem 7.1, the expected core-size formula for an irreducible rank $n$ Cartan type, restricted to the two non-simply-laced classical cases $X_n=B_n$ and $X_n=C_n$ with $b$ coprime to the Coxeter number. The statement text is copied exactly from its original authored declaration (native source0.tex lines 466-470), which the published article reprints in section 7 as Theorem 7.1 via its restatable macro; the environment wrapper here is editorial and the published counter is restored with setcounter. The complete author proof of the $B_n$ and $C_n$ cases is reproduced exactly from native lines 1278-1357, including both subsubsections 7.1.1 (Type $B$) and 7.1.2 (Type $C$). Retained as uncounted necessary core: the section 7 heading, the reduction of the expectation to a sum over $b\ac\cap\chk{\Q}$ (native 1216-1220), Proposition 7.2 (native 1222-1230), the coweight-lattice identity (7) (native 1234-1237), the Ehrhart quasipolynomial setup and Proposition 7.3 with its complete proof (native 1239-1253), the residue-class notation $\ac_{\sizeb}(b)_i$ (native 1255-1257) and the root-system data table Figure 5 (native 1259-1274). Sections 1-6, the simply-laced and exceptional cases, subsection 7.2 (types $F_4$ and $G_2$) and the rest of the paper are not selected. Original printed slips are preserved literally without repair, including the doubled word in native line 1281 and the mismatched index in native line 1355 where the type $C_n$ sum is written over $\cw_1$ although the surrounding computation indexes the points $k\cw_n$.
\setcounter{section}{6}
\section{Expected size of simultaneous cores}
\label{sec:expectation}

We are now ready to prove \Cref{thm:main_thm}:


\setcounter{theorem}{0}\begin{theorem}\label{thm:main_thm}
For $X_n$ an irreducible rank $n$ Cartan type with root system $\Phi$,
\begin{equation*}
\Expt{q \in \core(X_n,b)}{\chk{\size}(q)}=\frac{r \chk{g}}{h}\frac{n(b-1)(h+b+1)}{24},
\end{equation*}
where $h$ is the Coxeter number of $X_n$, $\chk{g}$ is the dual Coxeter number for $\chk{\Phi}$, and $r$ is the ratio of the length of a long root to the length of a short root in $\Phi$.

\end{theorem}
\begin{equation*}
    \Expt{q \in \core(X_n,b)}{\chk{\size}(q)} 
    = \frac{1}{|\core(X_n,b)|}\sum_{q\in \core(X_n,b)} \chk{\size}(q) = \frac{1}{|b\ac\cap\chk{\Q}|}\sum_{q\in b\ac\cap\chk{\Q}} \sizeb(q).
\end{equation*} 
The denominator was explicitly and nearly-uniformly calculated by Haiman \cite{haiman1994conjectures}. To compute the sum, we first record the vertices of the fundamental alcove $\ac$: they are \(\Gamma:=\{0\}\cup\left\{\frac{\cw_i}{c_i} : 1 \leq i \leq n\right\}\), where we recall that the $c_i$ are defined by $\widetilde{\alpha}=\sum_{i=1}^nc_i\alpha_i$.  As in \cite{thiel2017strange}, we proceed by translating the problem to the coweight lattice.  Define the \defn{extended affine Weyl group} by $\wex:=W\ltimes \cwl$, and write the group of automorphisms for $b\ac$ as $b\cycl:=\{\waf\in\wex:\waf(b\ac)=b\ac\}$. These groups are isomorphic for all $b$, and in particular have constant order that we denote by $f$.

\begin{prop}[{\cite[Theorem 2.5 \& Lemma 6.11]{thiel2017strange}}]
For any $b$ coprime to the Coxeter number $h$ in type $X_n$:
	\begin{enumerate}[(a)]
	\item The action of $b\cycl$ on $b\ac\cap\cwl$ is free.
	\item Each $b\cycl$ orbit of $b\ac\cap\cwl$ has exactly one element of $b\ac\cap\chk{\Q}=\core(X_n,b)$.
	\item For any $\waf\in b\cycl$ and any $\omega\in\cwl$, $\size^{(b)}(\omega)=\size^{(b)}(\waf\cdot\omega)$
	\end{enumerate}
\label{prop:barational}
\end{prop}
%REMARK: Part (c) isn't actually stated in the previous paper because of the incorrect definition of size, although the proof is identical--- it even uses $\wf=t_{\frac{b}{h}\chk{\rho}}wt_{-\frac{b}{h}\chk{\rho}}$.

Using this, we finish the translation from $\chk{\Q}$ to $\cwl$:

\setcounter{equation}{6}
\begin{equation}
\label{eq:ev_coweight_sum}
\Expt{q \in \core(X_n,b)}{\chk{\size}(q)} = \frac{1}{|b\ac\cap\chk{\Q}|}\cdot \frac{1}{f}\sum_{q\in b\ac\cap\cwl} \sizeb(q).    
\end{equation}

We now recall the relevant Ehrhart-theoretic tools. For any degree-$r$ polynomial $F:\R^n\to\R$, its \defn{weighted lattice point enumerator} over $\cwl$ is
$ \ac_F(b) := \sum\limits_{q\in b\ac\cap\cwl} F(q)$. This $\ac_F(b)$ is a quasipolynomial in $b$, of degree $n+r$ and period $c=c(\widetilde{X}_n):=\lcm(c_1,\dots,c_n)$, where the $c_i$ are again the denominators of the vertices of $\ac$. As $\size^{(b)}$ changes with $b$, Ehrhart theory appears to be inapplicable--- however, a judicious rewriting shows that this is not the case.

\begin{prop}
	The weighted lattice point enumerator $\ac_{\sizeb}(b)$
    is a quasipolynomial in $b$ of degree $n+2$ and period $c(\widetilde{X}_n):=\lcm(c_1,\ldots,c_n)$.
\label{prop:quasipoly}
\end{prop}

\begin{proof}
    Notice that $\chk{\size}_b(x)=\frac{h}{2}\left\|x\right\|^2-b\langle x,\chk{\rho}\rangle+(b^2-1)\frac{\left\|\chk{\rho}\right\|^2}{2h}.$
   Thus we find that
   \[\ac_{\sizeb}(b)={\textstyle\frac{h}{2}}\ac_{\|\cdot\|^2}(b)-b\ac_{\langle\cdot,\chk{\rho}\rangle}(b)+(b^2-1)\ac_{\frac{\left\|\chk{\rho}\right\|^2}{2h}}(b)\]
   is a quasipolynomial in $b$ of degree $n+2$ and period $c(\widetilde{X}_n)$.
\end{proof}

Therefore, to complete the proof of \Cref{thm:main_thm}, we may compute the quasipolynomial on for all components that contain a residue $b\mod c_{X_n}$ that is coprime to $h$. For the exceptional types, this is already a finite and computationally feasible calculation. 

Types $A$ and $D$ are simply-laced, and so their calculation was already completed in \cite{thiel2017strange}. Thus, we proceed along similar lines for types $B$ and $C$. In what follows, we write $\ac_{\sizeb}(b)_i$ to mean the polynomial which agrees with $\ac_{\sizeb}(b)$ for all $b\equiv i\bmod c(\widetilde{X}_n)$.  Relevant data to complete these computations for the irreducible root systems is provided in~\Cref{fig:data}.


\setcounter{figure}{4}
\begin{figure}[htbp]
\[\begin{array}{c|c|c|c|c|c|c|c}
X_n & h & \chk{g} & e_i & c_i & f & r & m(X_n) \\ \hline
A_n & n+1 & n+1 & 1,2,\ldots,n & 1,1,\ldots,1,1 & n+1 & 1 & 1 \\
B_n & 2n & n+1 & 1,3,\ldots,2n-1 & 1,2,\ldots,2,2 & 2 & 2 & 2\\
C_n & 2n & 2n-1 & 1,3,\ldots,2n-1 & 2,2,\ldots,2,1 & 2 & 2 & 2\\
D_n & 2n-2 & 2n-2 & 1,3,\ldots,2n-3,n-1 & 1,2,\ldots,2,1,1 & 4 & 1 & 2 \\
E_6 & 12 & 12 & 1, 4, 5, 7, 8, 11 & 1, 2, 2, 3, 2, 1 & 3 & 1 & 6 \\
E_7 & 18 & 18 & 1, 5, 7, 9, 11, 13, 17 & 2, 2, 3, 4, 3, 2, 1 & 2 & 1 & 12 \\
E_8 & 30 & 30 & 1, 7, 11, 13, 17, 19, 23, 29 & 2, 3, 4, 6, 5, 4, 3, 2 & 1 & 1 & 60 \\
F_4 & 12 & 9 & 1,5,7,11 & 2, 3, 4, 2 & 1 & 2 & 12 \\
G_2 & 6 & 4 & 1,5 & 3,2 & 1 & 3 & 6 \\
\end{array}\]
\caption{The type $X_n$, Coxeter number $h$, dual Coxeter number $\protect\chk{g}$, exponents $e_i$, coefficients of the highest root $c_i$, index of connection $f$, ratio of long to short root $r$, and $\mathrm{gcd}$ of the $c_i$ for the irreducible root systems.}
\label{fig:data}
\end{figure}

\subsection{Types B and C}

For either $X_n=B_n$ or $X_n=C_n$ we have $c(\widetilde{X}_n)=2$ and exponents $1,3,\ldots,2n-1$. Since all $b$ in the theorem statement are coprime to $h$, they must be odd, and so it suffices to compute only the polynomial $\ac_{\sizeb}(b)_1$. This polynomial has degree $n+2$. Recall the fact that the interior of $e_i\ac$ contains no points of $\cwl_{X_n}$ for all exponents $e_i$, because they are strictly less than $h$ (see e.g. \cite[Section 7.4]{thiel2017strange}). 
Thus, by Ehrhart reciprocity we conclude that $\ac_{\sizeb}(-e_i)=0$, and because in types $B_n$ and $C_n$ the exponents are the odd integers $1,3,\dots, 2n-1$, we thus we have $n$ roots of our desired polynomial. 

Moreover, it is easy to compute that $\ac_{\size}(1)=0$, and therefore (c.f. \cite[Proposition 7.5]{thiel2017strange}) we also have $\ac_{\sizeb}(-h-1)=0$. Hence we know all $n+2$ of our desired polynomial's roots, that is:
$$\ac_{\sizeb}(b)_1 = \kappa(b-1)(b+2n+1)\prod_{j=1}^{n} (b+e_j).$$
where the leading coefficient $\kappa$ depends only on $n$ and whether we are in types $B_n$ or $C_n$. Importantly, this formula holds for any odd $b$, not merely those $b$ which are coprime to $h=2n$. Thus we can determine it by explicitly calculating $\ac_{\sizeb}(b)$ at $b=3$, as follows:
\[ 
\kappa = \frac{\ac_{\size^{(3)}}(3)}{(3-1)(3+2n+1)\prod_{j=1}^{n} (3+(2j-1))} = \frac{1}{2^{n+2}(n+2)!}\ac_{\size^{(3)}}(3).
\]

From here, we recall Haiman's enumeration \cite{haiman1994conjectures}:  $|b\ac\cap\chk{\Q}|=\frac{1}{|W|}\prod_{j=1}^n(b+e_j)$, whenever $b$ and $h$ are coprime. Thus, we rewrite \Cref{eq:ev_coweight_sum} as the quadratic polynomial:
\begin{align*}
\Expt{q \in \core(X_n,b)}{\chk{\size}(q)} &= \frac{2^n n!\kappa}{f}(b-1)(b+2n+1) \\
&= \frac{2^n n!}{2\cdot 2^{n+2}(n+2)!}\ac_{\size^{(3)}}(3)\cdot (b-1)(b+h-1) \\
& = \ac_{\size^{(3)}}(3)\cdot \frac{(b-1)(b+h-1)}{8(n+1)(n+2)}
\end{align*}
Thus, to prove the formula from \Cref{thm:main_thm}, what remains to be verified is that
\[
\ac_{\size^{(3)}}(3)=\frac{r\chk{g}}{3h}\cdot n(n+1)(n+2) =
\begin{cases}
\frac{1}{3}(n+1)^2(n+2) & \text{in type } B_n, \\
\frac{1}{3}(n+1)(n+2)(2n-1) & \text{in type } C_n.
\end{cases}
\]


%------%
\subsubsection{Type B}
\label{sec:type_b_calc}

In type $B_n$, since we have $\sum_{j=1}^n a_j \cw_j = (\sum_{j=1}^n a_j,\sum_{i=2}^n a_j,\ldots,a_n)$ and $\widetilde{\alpha}=(1,1,0,0,\ldots,0)$, the coweight points in $3\ac$ are those points for which $a_1+2(\sum_{j=2}^n a_j) \leq 3$, which are the $2n+2$ points \[\{0,\cw_1,2\cw_1,3\cw_1\}\cup \{\cw_j,\cw_1+\cw_j\}_{j=2}^n.\]
We compute that for $2 \leq j \leq n$
\begin{align*}\size^{(3)}&(\cw_j) = \frac{n}{(2n)^2}\left( \left(\|2n\cw_j-3\chk{\rho}\|\right)^2-\|\chk{\rho}\|^2\right)
\\ &=
\frac{1}{4n} \left( \sum_{i=1}^j \big(2n-3(n+1-i)\big)^2+\sum_{i=j+1}^n \big(3(n+1-i)\big)^2 -\frac{n(n+1)(2n+1)}{6}\right) \\ 
 &= \frac{(2n+2-3 j) (2n+1-3 j)}{6}. 
\end{align*}
A similar computation also shows that $\size^{(3)}(\cw_1+\cw_j) = \frac{(2n+2-3 j) (2n+1-3 j)}{6}$
and therefore 
\[\sum_{j=2}^n \size^{(3)}(\cw_j)+\size^{(3)}(\cw_1+\cw_j) = \frac{(n-2) (n-1)^2}{3}.\]
The remaining points give contributions of the following form: 
\begin{align*}
\size^{(3)}(k\cw_1) &= \frac{n}{(2 n)^2} \left( (2 k n - 3n)^2 + 
     \sum_{i=2}^n (3 (n + 1 - i))^2 - \sum_{i=1}^n (n + 1 - i)^2 \right) \\
    &=\frac{1 + 3 n - 9 k n + 3 k^2 n + 2 n^2}{3}.
\end{align*}

We thus compute that $\sum_{k=0}^3 \size^{(3)}(k\cw_1)= \frac{4(1 + 2 n^2)}{3}$. Putting these together, we obtain 
$$ \ac_{\size^{(3)}}(3) = \frac{{n-2}{n-1}^2}{3} + \frac{4(1+2n^2)}{3} = \frac{1}{3}(n+1)^2(n+2) $$
which concludes the proof of \Cref{thm:main_thm} for type $B_n$.


%------%
\subsubsection{Type C}
\label{sec:type_c_calc}
Turning to type $C_n$, we have $\widetilde{\alpha}=\frac{2}{\sqrt{2}}e_1$ and
$$\sum_{j=1}^n a_j \cw_j = \sqrt{2}\left(\frac{a_n}{2}+\sum_{j=1}^{n-1} a_j,~\frac{a_n}{2}+\sum_{j=2}^n a_j,\ldots,~\frac{a_n}{2}+a_{n-1},~\frac{a_n}{2}\right)$$
So the coweight points in $3\ac$ are those points for which $2(\sum_{j=1}^{n-1} a_j)+a_n \leq 3$, which are the $2n+2$ points

\[\{0,\cw_n,2\cw_n,3\cw_n\}\cup \{\cw_j,\cw_n+\cw_j\}_{j=1}^{n-1}.\]
Using the fact that $\chk{\rho}=\sum_{j=1}^n \cw_j = \frac{\sqrt{2}}{2}(2n-1,2n-3,2n-5,\ldots,1)$, we compute that for $1 \leq j \leq n-1$
\begin{align*}
\size^{(3)}&(\cw_j) = \frac{n}{(2n)^2}\left( \left(\|2n\cw_j-3\chk{\rho}\|\right)^2-\|\chk{\rho}\|^2\right) \\ &=
\frac{1}{4n}\left( \sum_{i=1}^j 2\big(2n-\frac{3}{2}(2n+1-2i)\big)^2+\sum_{i=j+1}^n \frac{(3(2n+1-2i)\big)^2}{2} -\frac{4n^3-n}{6}\right) \\ 
 &= \frac{(2n-3j)^2-1}{3}.
\end{align*}
A similar computation also shows that $\size^{(3)}(\cw_n+\cw_j) = \frac{(n-3j)^2-1}{3}$ and therefore
\[\sum_{j=1}^{n-1} \size^{(3)}(\cw_j)+\size^{(3)}(\cw_n+\cw_j) = \frac{(n-2) (n-1)(2n+1)}{3}.\]
The remaining points give contributions of the following form:
\begin{align*}
\size^{(3)}(k\cw_n) &= \frac{n}{(2 n)^2} \left(
     \sum_{i=1}^n 2 \left(kn-\frac{3}{2}(2n+1-2i)\right)^2 -\sum_{i=1}^n \frac{(2n+1-i)^2}{2} \right) \\
    &=\frac{-2+8 n^2-9 k n^2+3 k^2 n^2}{6}.
\end{align*}

We thus compute that $\sum_{k=0}^3 \size^{(3)}(k\cw_1)  = \frac{2(5n^2-2)}{3}$. Putting these together, we obtain
$$ \ac_{\sizeb}(b) = \frac{1}{3}(n+1)(n+2)(2n-1),$$
which concludes the proof of \Cref{thm:main_thm} for type $C_n$.

\begin{thebibliography}{99}
\bibitem[10]{haiman1994conjectures} Mark D. Haiman, \emph{Conjectures on the quotient ring by diagonal invariants}, J. Algebraic Combin. \textbf{3} (1994), no. 1, 17--76.
\bibitem[20]{thiel2017strange} Marko Thiel and Nathan Williams, \emph{Strange expectations and simultaneous cores}, J. Algebraic Combin. \textbf{46} (2017), no. 1, 219--261.
\end{thebibliography}
\end{document}
